\documentclass[11pt,a4paper]{article}
\usepackage{booktabs}

\begin{document}

\section{Tables}
Table~\ref{tab:simple} has vertical and horizontal rules, and
Table~\ref{tab:booktabs} is set in the style of the \texttt{booktabs}
package, with fewer and lighter rules.

\begin{table}[h]
\centering
\caption{A simple table.}
\label{tab:simple}
\begin{tabular}{|l|c|r|}
\hline
Name & Count & Price \\
\hline
Apples & 3 & 1.20 \\
Pears & 12 & 0.85 \\
\hline
\multicolumn{2}{|l|}{Total} & 2.05 \\
\hline
\end{tabular}
\end{table}

Text continues after the first table and runs long enough to wrap onto a
second line, so the spacing after a float can be measured.

\begin{table}[h]
\centering
\caption{Energy return on investment by fuel: a caption long enough to need
two lines in the article class, which then sets it as a paragraph.}
\label{tab:booktabs}
\begin{tabular}{@{}lrr@{}}
\toprule
Fuel & EROI & Year \\
\midrule
Coal & 46 & 2010 \\
Oil and gas & 16 & 2010 \\
\bottomrule
\end{tabular}
\end{table}

After both tables the text continues.

\end{document}
